\subsection{Defect Detection Enhancement}
\label{sec:results-detection}

Defect detection performance was evaluated by applying a standard arrival-time-based defect localization algorithm to denoised waveforms. The algorithm computes the time-of-arrival (TOA) difference between two ultrasonic transducers and uses the resulting ellipses to triangulate defect positions; a threshold on the ellipse intersection confidence determines detection. Classification metrics were computed at three averaging levels (16, 64, and 256 averages) across all defect-containing specimens.

\textbf{Classification metrics.} At 16 averages, PMDF-Net achieved an F1-score of \placeholder{F1_pmdf_16} compared to \placeholder{F1_baseline_16} for the best conventional baseline (Wiener filtering), representing a relative improvement of \placeholder{F1_rel_improvement}\%. The ROC AUC for PMDF-Net at this averaging level was \placeholder{AUC_pmdf_16} versus \placeholder{AUC_baseline_16} for Wiener filtering. As the number of averages increased, all methods improved, but PMDF-Net maintained a consistent lead. At 256 averages, PMDF-Net's F1 reached \placeholder{F1_pmdf_256} while the best baseline (BM3D) achieved \placeholder{F1_bm3d_256}, indicating that PMDF-Net with 16 averages approximates the detection performance of conventional methods with 256 averages.

\textbf{POD curves.} Figure~\ref{fig:detection_results} presents the probability of detection (POD) as a function of defect size. PMDF-Net achieved 90\% POD at a minimum defect size of \placeholder{POD90_pmdf} mm, compared to \placeholder{POD90_baseline} mm for the best baseline. The POD curve for PMDF-Net rises more steeply, indicating reliable detection of smaller defects. The minimum detectable defect size at 80\% POD was \placeholder{MDS80_pmdf} mm for PMDF-Net versus \placeholder{MDS80_baseline} mm for Wiener filtering.

\textbf{Efficiency gain.} The crossover between PMDF-Net at reduced averaging and conventional methods at full averaging quantifies the practical efficiency benefit. PMDF-Net operating at 16 averages demonstrated detection performance comparable to conventional signal averaging at 256 averages, corresponding to a \textbf{16$\times$} reduction in acquisition time. At equivalent detection performance (F1 = \placeholder{F1_equiv}), PMDF-Net required only \placeholder{avg_pmdf} averages while Wiener filtering required \placeholder{avg_wiener} averages, confirming the efficiency advantage across the full defect size range.

\textbf{Feature preservation driving detection improvement.} The gains in defect detection arise directly from PMDF-Net's superior feature preservation. Accurate arrival time estimates (Section~\ref{sec:results-imaging}) reduce localization error in TOA-based triangulation, producing tighter position confidence ellipses and fewer missed detections at small defect sizes. Preserved peak amplitudes maintain the signal-to-noise ratio of the arrival envelope, enabling reliable defect sizing even at low averaging levels. These results confirm that waveform fidelity, not just aggregate SNR, is the operative factor for practical defect detection.
